% Nouveau EGA V (cote 25): archive cote 25, site page 29 (PDF page 30).
% typescript. Draft compared: none.
% AI-transcribed by Claude Opus 5.5 (claude-opus-5-5, adaptive thinking, effort high), Message Batches API, batch ega25b (msgbatch_01A7n5YuXubUPNK1zsepi8RG), 2026-11-13; not yet checked by eye.
\scanpage{25/29}
\mspage{\hand{25}}

Je dois cependant d'abord signqler un remords à 4.2. et 4.3. , qui
dev-aient être énoncés sous une forme légèrement plus générale (tout au
moins en remarque): on se donne un entier m, et on désigne par E l'ensem
ble des \hand{$\xi\in$}$\overset{\hand{\vee}}{\text{P}}$ tels que G$_{\hand{\xi}}$ resp. Y$_{\hand{\xi}}$ satisfasse P sauf \var{\deleted{au plus}}{}{deletion}{ega25_25_029-v0001} sur une
partie fermée de dimension\hand{$\leq$}m (i.e. l'ensemble P-singulier Z est de dimen-%
sion\hand{$\leq$} m). Alors a) E est une partie constructible de $\overset{\hand{\vee}}{\text{P}}$ et b) dans le
cas S=Spec(k), si F \var{\inserted[typed]{resp. X}}{resp. X}{insertion}{ega25_25_029-v0002} satisfait à P sauf \var{\deleted{au plus}}{}{deletion}{ega25_25_029-v0003} sur un ensemble de
\var{\xout{xx}}{}{typewriter strike-out}{ega25_25_029-v0004}dimension\hand{$\leq$} m+1 , alors E contient un ouvert non vide.

\uline{Proposition} 7.2. Supposons que X soit propre sur S, et F de présentati-%
on finie et plat sur S. Soit P l'une des propriétés (i) à (v) de 4.2.
et soit m un entier. Soit s\hand{$\in$}S, et supposons que \var{\xout{k(s)xxF$_s$xsatisfassexà}}{}{typewriter strike-out}{ega25_25_029-v0005}
\var{\xout{laxpropriétéxP}}{}{typewriter strike-out}{ega25_25_029-v0006} l\var{\xout{'}}{}{typewriter strike-out}{ega25_25_029-v0007}ensemble Z$_s$ des points de X$_s$ en lesquels F$_s$ ne satisfait
pas à P soit de dimension \hand{$\leq$} m+1. Alors, si de plus k(s) est infini, il
existe un voisinage U de s dans E, et une section \hand{$\xi$} de $\overset{\hand{\vee}}{\text{P}}$ sur U, ayant
les propriétés suivantes: pour tout s\hand{$\in$}U, l'ensemble Z des points
de \var{\deleted{X}}{}{deletion}{ega25_25_029-v0008} Y$_{\hand{\xi,\xi(\textcolor{black}{s})}}$ en lesquels G$_{\hand{\xi(\textcolor{black}{s})}}$ ne satisfait pas à P est de dimension
\hand{$\leq$}m, et $\phi_{\hand{\xi}(s)}$ est F$_{\hand{\xi,\xi(\textcolor{black}{s})}}$-régulière\hand{,} \var{\xout{xlexModulexGxxsur}}{}{typewriter strike-out}{ega25_25_029-v0009} Sous ces condi-%
tions, le Module G$_{\hand{\xi}}$ sur Y$_{\hand{\xi}}$ est plat relativement à U. Enfin, si k(s)
n'est pas supposé inffni, on peut encore faire la construction précéden-%
te, après extension étale duntype envisagé dans 7.1.

\uline{Proposition} 7.3. Essentiellement la même, il n'y a plus de F, \unsure{on|ok}
\var{\margin{[N.B. Pour (v) cf 12.2.1. (x) et (xi) On peut donc ajouter \unsure{un|une} autre cas à 4.3. par 12.2.1. \unsure{(x)}]}}{\placedmargin{[N.B. Pour (v) cf 12.2.1. (x) et (xi) On peut donc ajouter \unsure{un|une} autre cas à 4.3. par 12.2.1. \unsure{(x)}]}}{marginal note placed}{ega25_25_029-v0010}
suppose X plat rel à S, on réfère aux propriétés (i) à \var{\deleted{(xxx)}}{}{deletion}{ega25_25_029-v0011}\var{\inserted[typed]{(\unsure{\var{\deleted{i}}{}{deletion}{ega25_25_029-v0013}v|iv})}}{(\unsure{\var{\deleted{i}}{}{deletion}{ega25_25_029-v0013}v|iv})}{insertion}{ega25_25_029-v0012} de 4.3. au
lieu de celles de 4.2. (mais en ayant soin de faire \var{\deleted{les réserves}}{}{deletion}{ega25_25_029-v0014} \var{\xout{axxxxxxxxxxx}}{}{typewriter strike-out}{ega25_25_029-v0015}\var{\inserted{la réserve}}{la réserve}{insertion}{ega25_25_029-v0016}
``k\var{\inserted[typed]{(s)}}{(s)}{insertion}{ega25_25_029-v0017} de car. 0 ou f:X\hand{$\longrightarrow$}P une immersion\hand{\unsure{$\xi$}}\var{\inserted{, et dans le cas (v), pour tout $s\in S$, et tte comp. irréductible Z de X$_s$, on ait $\dim f(Z)\geq 2$).}}{, et dans le cas (v), pour tout $s\in S$, et tte comp. irréductible Z de X$_s$, on ait $\dim f(Z)\geq 2$).}{insertion}{ega25_25_029-v0018}

\var{\deleted{On notera que nous avons exclus les propriétés (v) et (vi) de}}{}{deletion}{ega25_25_029-v0019}
\var{\deleted{4.3. , qui ne sont pas de nature purement locale. Cependant, on récupère}}{}{deletion}{ega25_25_029-v0020}
\var{\deleted{le cas (v) en notant qu'on peut s'arranger en vertu de la démonstratio}}{}{deletion}{ega25_25_029-v0021}
\var{\deleted{de 7.3. que Y$_{\hand{\xi}}$ soit plat sur U à fibres séparables, mais alors en}}{}{deletion}{ega25_25_029-v0022}
\var{\margin{29}}{\placedmargin{29}}{marginal note placed}{ega25_25_029-v0023}
